
\documentclass{article}
\usepackage{colortbl}
\usepackage{makecell}
\usepackage{multirow}
\usepackage{supertabular}

\begin{document}

\newcounter{utterance}

\centering \large Interaction Transcript for game `clean\_up', experiment `1\_medium\_3obj\_ar', episode 4 with Qwen3.6{-}35B{-}A3B{-}FP8{-}without{-}reasoning.
\vspace{24pt}

{ \footnotesize  \setcounter{utterance}{1}
\setlength{\tabcolsep}{0pt}
\begin{supertabular}{c@{$\;$}|p{.15\linewidth}@{}p{.15\linewidth}p{.15\linewidth}p{.15\linewidth}p{.15\linewidth}p{.15\linewidth}}
    \# & $\;$A & \multicolumn{4}{c}{Game Master} & $\;\:$B\\
    \hline

    \theutterance \stepcounter{utterance}  
    & & \multicolumn{4}{p{0.6\linewidth}}{
        \cellcolor[rgb]{0.9,0.9,0.9}{
            \makecell[{{p{\linewidth}}}]{
                \texttt{\tiny{[P1$\langle$GM]}}
                \texttt{أنا مدير لعبتك، وأنت تلعب لعبة تعاونية مع الشبكة التالية كلوحة لعب:} \\
\\ 
\texttt{```} \\
\texttt{ 1234567} \\
\texttt{╔═╤══╤══╗} \\
\texttt{║◌│◌◌│◌◌║ 1} \\
\texttt{╟─┤◌◌│◌◌║ 2} \\
\texttt{║◌│◌◌├──╢ 3} \\
\texttt{╟─┘C◌│◌◌║ 4} \\
\texttt{║◌◌◌◌│◌◌║ 5} \\
\texttt{║◌L┌─┼──╢ 6} \\
\texttt{║◌◌│◌│P◌║ 7} \\
\texttt{╚══╧═╧══╝} \\
\\ 
\texttt{```} \\
\\ 
\texttt{* تعرض الحافة العلوية إحداثيات x التي تتزايد نحو اليمين، والحافة اليمنى إحداثيات y التي تتزايد نحو الأسفل.} \\
\texttt{* الكائنات التالية موضوعة عشوائياً على شبكتك: 'C', 'L', 'P'.} \\
\\ 
\\ 
\texttt{يرى اللاعب الآخر نسخة مختلفة من لوحة اللعب، حيث توضع الكائنات في مواقع عشوائية مختلفة. لا يمكنك رؤية لوحة اللاعب الآخر، ولا يمكنه رؤية لوحتك.} \\
\\ 
\\ 
\texttt{**الهدف:**} \\
\\ 
\texttt{يجب على كلا اللاعبين تحريك الكائنات على خلفية كل منهما بحيث تنتهي الكائنات المتطابقة عند الإحداثيات نفسها. عليك التواصل مع اللاعب الآخر للاتفاق على إعداد هدف مشترك.} \\
\\ 
\\ 
\texttt{**القواعد:**} \\
\\ 
\texttt{* في كل دور، يمكنك إرسال أمر واحد فقط من الأمرين التاليين:} \\
\texttt{ 1. `قل: <رسالة>`: لإرسال رسالة (كل شيء حتى فاصل السطر التالي) إلى اللاعب الآخر. سأقوم بتمريرها إلى شريكك.} \\
\texttt{2. `حرّك: <كائن>, (<X>, <Y>)`: لتحريك كائن إلى موضع جديد، حيث `<X>` هو العمود و `<Y>` هو الصف. سأخبرك بما إذا كانت حركتك ناجحة أم لا.} \\
\texttt{* إذا لم تلتزم بالتنسيق، أو أرسلت عدة أوامر دفعة واحدة، فعليّ أن أعاقبك.} \\
\texttt{* إذا تجاوز كلا اللاعبين أكثر من 8 عقوبة، فستخسران اللعبة معاً.} \\
\texttt{* من الضروري أن تتواصل مع اللاعب الآخر بشأن حالة هدفك! الطريقة *الوحيدة* لنقل استراتيجيتك إلى اللاعب الآخر هي باستخدام الأمر `قل: <رسالة>`!} \\
\\ 
\\ 
\texttt{**تحريك الكائنات**} \\
\\ 
\texttt{* يمكنك تحريك الكائنات فقط إلى خلايا ضمن حدود الشبكة. يجب أن تكون الخلية المستهدفة فارغة، أي يجب أن تحتوي فقط على الرمز '◌'.} \\
\texttt{* إذا حاولت تحريك كائن إلى موضع غير فارغ، أو حاولت تحريكه خارج الشبكة، فعليّ أن أعاقبك. ستحصل على محاولة أخرى.} \\
\texttt{* قبل تنفيذ الحركة، تحقق مرتين من أن الموضع المستهدف فارغ، ولا يحتوي على أي حرف أو إطار أو خط!} \\
\\ 
\\ 
\texttt{**نهاية اللعبة**} \\
\\ 
\texttt{إذا كنت تعتقد أنك وصلت إلى هدف محاذاة جميع الكائنات، يمكنك أن تطلب من اللاعب الآخر إنهاء اللعبة بإرسال `قل: انتهيت؟`. إذا طلب منك اللاعب الآخر إنهاء اللعبة، ورددت بـ `قل: انتهيت!`، ستنتهي اللعبة.} \\
\\ 
\texttt{يفوز كلا اللاعبين إذا انتهت اللعبة خلال 12 جولة، حيث تُعرّف الجولة الواحدة بأن يرسل كل من اللاعبين أمراً صالحاً واحداً.} \\
\\ 
\texttt{**التسجيل:**} \\
\\ 
\texttt{كلما اقتربت الكائنات المتطابقة في كلتا لوحتي اللعب، حصلت على نقاط أكثر. العقوبات تقلل نقاطك. هل يمكنك تحطيم الرقم القياسي؟} \\
\\ 
\\ 
\texttt{يرجى إرسال رسالة إلى اللاعب الآخر لبدء اللعبة!} \\
            }
        }
    }
    & \\ \\

    \theutterance \stepcounter{utterance}  
    & \multicolumn{4}{p{0.6\linewidth}}{
        \cellcolor[rgb]{0.9,0.9,0.9}{
            \makecell[{{p{\linewidth}}}]{
                \texttt{\tiny{[P1$\rangle$GM]}}
                \texttt{قل: مرحباً! لقد راجعت اللوحة الخاصة بي. أرى الكائنات C و L و P في مواقع مختلفة. بما أننا لا نرى نفس التوزيع، يجب أن نتفق على استراتيجية موحدة.} \\
\\ 
\texttt{لنبدأ بتحديد هدف مشترك. هل توافق على أن نحاول تجميع جميع الكائنات الثلاثة (C, L, P) في الزاوية العلوية اليسرى (الإحداثيات 1,1)؟ أو ربما في المركز؟} \\
\\ 
\texttt{أقترح أن نبدأ بالحركة نحو الزاوية (1,1) لأنها نقطة مرجعية سهلة. هل أنت موافق على محاولة تجميعهم هناك؟ إذا كنت توافق، سأبدأ بتحريك أقرب كائن. يرجى تأكيد استراتيجيتك أو اقتراح إحداثيات هدف مختلفة.} \\
            }
        }
    }
    & & \\ \\

    \theutterance \stepcounter{utterance}  
    & & & \multicolumn{2}{p{0.3\linewidth}}{
        \cellcolor[rgb]{0.9,0.9,0.9}{
            \makecell[{{p{\linewidth}}}]{
                \texttt{\tiny{[GM$|$GM]}}
                \texttt{قل: مرحباً! لقد راجعت اللوحة الخاصة بي. أرى الكائنات C و L و P في مواقع مختلفة. بما أننا لا نرى نفس التوزيع، يجب أن نتفق على استراتيجية موحدة.} \\
\\ 
\texttt{لنبدأ بتحديد هدف مشترك. هل توافق على أن نحاول تجميع جميع الكائنات الثلاثة (C, L, P) في الزاوية العلوية اليسرى (الإحداثيات 1,1)؟ أو ربما في المركز؟} \\
\\ 
\texttt{أقترح أن نبدأ بالحركة نحو الزاوية (1,1) لأنها نقطة مرجعية سهلة. هل أنت موافق على محاولة تجميعهم هناك؟ إذا كنت توافق، سأبدأ بتحريك أقرب كائن. يرجى تأكيد استراتيجيتك أو اقتراح إحداثيات هدف مختلفة.} \\
            }
        }
    }
    & & \\ \\

    \theutterance \stepcounter{utterance}  
    & & & \multicolumn{2}{p{0.3\linewidth}}{
        \cellcolor[rgb]{0.9,0.9,0.9}{
            \makecell[{{p{\linewidth}}}]{
                \texttt{\tiny{[GM$|$GM]}}
                \texttt{Invalid format: tail is not empty: '} \\
\\ 
\texttt{لنبدأ بتحديد هدف مشترك. هل توافق على أن نحاول تجميع جميع الكائنات الثلاثة (C, L, P) في الزاوية العلوية اليسرى (الإحداثيات 1,1)؟ أو ربما في المركز؟} \\
\\ 
\texttt{أقترح أن نبدأ بالحركة نحو الزاوية (1,1) لأنها نقطة مرجعية سهلة. هل أنت موافق على محاولة تجميعهم هناك؟ إذا كنت توافق، سأبدأ بتحريك أقرب كائن. يرجى تأكيد استراتيجيتك أو اقتراح إحداثيات هدف مختلفة.} \\
            }
        }
    }
    & & \\ \\

    \theutterance \stepcounter{utterance}  
    & & \multicolumn{4}{p{0.6\linewidth}}{
        \cellcolor[rgb]{0.9,0.9,0.9}{
            \makecell[{{p{\linewidth}}}]{
                \texttt{\tiny{[P1$\langle$GM]}}
                \texttt{عقوبة: يجب ألا تحتوي رسالتك على أي شيء بعد الأمر!} \\
\\ 
\texttt{تأكد من أن ردك يحتوي فقط على `قل: <رسالة>` أو `حرّك: <كائن>, (<X>, <Y>)`، ولا شيء غير ذلك!} \\
\\ 
\texttt{لقد جمعتما معاً 1 من أصل 8 عقوبة.} \\
\texttt{يرجى المحاولة مرة أخرى!} \\
            }
        }
    }
    & \\ \\

    \theutterance \stepcounter{utterance}  
    & \multicolumn{4}{p{0.6\linewidth}}{
        \cellcolor[rgb]{0.9,0.9,0.9}{
            \makecell[{{p{\linewidth}}}]{
                \texttt{\tiny{[P1$\rangle$GM]}}
                \texttt{قل: مرحباً! أرى C في (4,4)، L في (6,6)، و P في (7,7). لنحاول تجميعهم في الزاوية العلوية اليسرى (1,1). هل توافق على هذه الاستراتيجية؟} \\
            }
        }
    }
    & & \\ \\

    \theutterance \stepcounter{utterance}  
    & & & \multicolumn{2}{p{0.3\linewidth}}{
        \cellcolor[rgb]{0.9,0.9,0.9}{
            \makecell[{{p{\linewidth}}}]{
                \texttt{\tiny{[GM$|$GM]}}
                \texttt{قل: مرحباً! أرى C في (4,4)، L في (6,6)، و P في (7,7). لنحاول تجميعهم في الزاوية العلوية اليسرى (1,1). هل توافق على هذه الاستراتيجية؟} \\
            }
        }
    }
    & & \\ \\

    \theutterance \stepcounter{utterance}  
    & & \multicolumn{4}{p{0.6\linewidth}}{
        \cellcolor[rgb]{0.9,0.9,0.9}{
            \makecell[{{p{\linewidth}}}]{
                \texttt{\tiny{[GM$\rangle$P2]}}
                \texttt{أنا مدير لعبتك، وأنت تلعب لعبة تعاونية مع الشبكة التالية كلوحة لعب:} \\
\\ 
\texttt{```} \\
\texttt{ 1234567} \\
\texttt{╔═╤══╤══╗} \\
\texttt{║◌│◌◌│◌◌║ 1} \\
\texttt{╟─┤◌◌│◌◌║ 2} \\
\texttt{║◌│L◌├──╢ 3} \\
\texttt{╟─┘◌◌│◌C║ 4} \\
\texttt{║P◌◌◌│◌◌║ 5} \\
\texttt{║◌◌┌─┼──╢ 6} \\
\texttt{║◌◌│◌│◌◌║ 7} \\
\texttt{╚══╧═╧══╝} \\
\\ 
\texttt{```} \\
\\ 
\texttt{* تعرض الحافة العلوية إحداثيات x التي تتزايد نحو اليمين، والحافة اليمنى إحداثيات y التي تتزايد نحو الأسفل.} \\
\texttt{* الكائنات التالية موضوعة عشوائياً على شبكتك: 'C', 'L', 'P'.} \\
\\ 
\\ 
\texttt{يرى اللاعب الآخر نسخة مختلفة من لوحة اللعب، حيث توضع الكائنات في مواقع عشوائية مختلفة. لا يمكنك رؤية لوحة اللاعب الآخر، ولا يمكنه رؤية لوحتك.} \\
\\ 
\\ 
\texttt{**الهدف:**} \\
\\ 
\texttt{يجب على كلا اللاعبين تحريك الكائنات على خلفية كل منهما بحيث تنتهي الكائنات المتطابقة عند الإحداثيات نفسها. عليك التواصل مع اللاعب الآخر للاتفاق على إعداد هدف مشترك.} \\
\\ 
\\ 
\texttt{**القواعد:**} \\
\\ 
\texttt{* في كل دور، يمكنك إرسال أمر واحد فقط من الأمرين التاليين:} \\
\texttt{ 1. `قل: <رسالة>`: لإرسال رسالة (كل شيء حتى فاصل السطر التالي) إلى اللاعب الآخر. سأقوم بتمريرها إلى شريكك.} \\
\texttt{2. `حرّك: <كائن>, (<X>, <Y>)`: لتحريك كائن إلى موضع جديد، حيث `<X>` هو العمود و `<Y>` هو الصف. سأخبرك بما إذا كانت حركتك ناجحة أم لا.} \\
\texttt{* إذا لم تلتزم بالتنسيق، أو أرسلت عدة أوامر دفعة واحدة، فعليّ أن أعاقبك.} \\
\texttt{* إذا تجاوز كلا اللاعبين أكثر من 8 عقوبة، فستخسران اللعبة معاً.} \\
\texttt{* من الضروري أن تتواصل مع اللاعب الآخر بشأن حالة هدفك! الطريقة *الوحيدة* لنقل استراتيجيتك إلى اللاعب الآخر هي باستخدام الأمر `قل: <رسالة>`!} \\
\\ 
\\ 
\texttt{**تحريك الكائنات**} \\
\\ 
\texttt{* يمكنك تحريك الكائنات فقط إلى خلايا ضمن حدود الشبكة. يجب أن تكون الخلية المستهدفة فارغة، أي يجب أن تحتوي فقط على الرمز '◌'.} \\
\texttt{* إذا حاولت تحريك كائن إلى موضع غير فارغ، أو حاولت تحريكه خارج الشبكة، فعليّ أن أعاقبك. ستحصل على محاولة أخرى.} \\
\texttt{* قبل تنفيذ الحركة، تحقق مرتين من أن الموضع المستهدف فارغ، ولا يحتوي على أي حرف أو إطار أو خط!} \\
\\ 
\\ 
\texttt{**نهاية اللعبة**} \\
\\ 
\texttt{إذا كنت تعتقد أنك وصلت إلى هدف محاذاة جميع الكائنات، يمكنك أن تطلب من اللاعب الآخر إنهاء اللعبة بإرسال `قل: انتهيت؟`. إذا طلب منك اللاعب الآخر إنهاء اللعبة، ورددت بـ `قل: انتهيت!`، ستنتهي اللعبة.} \\
\\ 
\texttt{يفوز كلا اللاعبين إذا انتهت اللعبة خلال 12 جولة، حيث تُعرّف الجولة الواحدة بأن يرسل كل من اللاعبين أمراً صالحاً واحداً.} \\
\\ 
\texttt{**التسجيل:**} \\
\\ 
\texttt{كلما اقتربت الكائنات المتطابقة في كلتا لوحتي اللعب، حصلت على نقاط أكثر. العقوبات تقلل نقاطك. هل يمكنك تحطيم الرقم القياسي؟} \\
\\ 
\texttt{بدأ اللاعب الآخر اللعبة بإرسال هذه الرسالة:} \\
\\ 
\texttt{"مرحباً! أرى C في (4,4)، L في (6,6)، و P في (7,7). لنحاول تجميعهم في الزاوية العلوية اليسرى (1,1). هل توافق على هذه الاستراتيجية؟"} \\
\\ 
\texttt{ما هو أمرك الأول؟} \\
            }
        }
    }
    & \\ \\

    \theutterance \stepcounter{utterance}  
    & & & \multicolumn{4}{p{0.6\linewidth}}{
        \cellcolor[rgb]{0.9,0.9,0.9}{
            \makecell[{{p{\linewidth}}}]{
                \texttt{\tiny{[GM$\langle$P2]}}
                \texttt{قل: مرحباً! أرى L في (3,3)، C في (4,4)، و P في (5,1). لا أوافق على تجميعهم في (1,1) لأن P بعيد جداً. لنحاول تجميعهم في الزاوية السفلية اليمنى (7,7) حيث C و L قريبان نسبياً، وسنحاول سحب P إليها. هل توافق؟} \\
            }
        }
    }
     \\ \\

    \theutterance \stepcounter{utterance}  
    & & & \multicolumn{2}{p{0.3\linewidth}}{
        \cellcolor[rgb]{0.9,0.9,0.9}{
            \makecell[{{p{\linewidth}}}]{
                \texttt{\tiny{[GM$|$GM]}}
                \texttt{قل: مرحباً! أرى L في (3,3)، C في (4,4)، و P في (5,1). لا أوافق على تجميعهم في (1,1) لأن P بعيد جداً. لنحاول تجميعهم في الزاوية السفلية اليمنى (7,7) حيث C و L قريبان نسبياً، وسنحاول سحب P إليها. هل توافق؟} \\
            }
        }
    }
    & & \\ \\

    \theutterance \stepcounter{utterance}  
    & & \multicolumn{4}{p{0.6\linewidth}}{
        \cellcolor[rgb]{0.9,0.9,0.9}{
            \makecell[{{p{\linewidth}}}]{
                \texttt{\tiny{[P1$\langle$GM]}}
                \texttt{تم تمرير رسالتك إلى اللاعب الآخر.} \\
\\ 
\\ 
\texttt{أنت تلعب حالياً الجولة 1 من أصل 12 كحد أقصى.} \\
\texttt{لقد جمعتما معاً 1 من أصل 8 عقوبة.} \\
\texttt{أرسل اللاعب الآخر هذه الرسالة:} \\
\\ 
\\ 
\texttt{"مرحباً! أرى L في (3,3)، C في (4,4)، و P في (5,1). لا أوافق على تجميعهم في (1,1) لأن P بعيد جداً. لنحاول تجميعهم في الزاوية السفلية اليمنى (7,7) حيث C و L قريبان نسبياً، وسنحاول سحب P إليها. هل توافق؟"} \\
\\ 
\\ 
\texttt{ما هو أمرك التالي؟} \\
            }
        }
    }
    & \\ \\

    \theutterance \stepcounter{utterance}  
    & \multicolumn{4}{p{0.6\linewidth}}{
        \cellcolor[rgb]{0.9,0.9,0.9}{
            \makecell[{{p{\linewidth}}}]{
                \texttt{\tiny{[P1$\rangle$GM]}}
                \texttt{قل: مقترح جيد. الزاوية (7,7) تبدو معقولة. سأبدأ بتحريك C من (4,4) إلى (7,7) إذا كانت الخلية فارغة. هل أنت متأكد من أن (7,7) فارغة حالياً؟ إذا كانت P لا تزال هناك، فلن أتمكن من التحرك إليها. يرجى تأكيد حالة الخلية (7,7) أو اقتراح إحداثيات بديلة إذا لزم الأمر.} \\
            }
        }
    }
    & & \\ \\

    \theutterance \stepcounter{utterance}  
    & & & \multicolumn{2}{p{0.3\linewidth}}{
        \cellcolor[rgb]{0.9,0.9,0.9}{
            \makecell[{{p{\linewidth}}}]{
                \texttt{\tiny{[GM$|$GM]}}
                \texttt{قل: مقترح جيد. الزاوية (7,7) تبدو معقولة. سأبدأ بتحريك C من (4,4) إلى (7,7) إذا كانت الخلية فارغة. هل أنت متأكد من أن (7,7) فارغة حالياً؟ إذا كانت P لا تزال هناك، فلن أتمكن من التحرك إليها. يرجى تأكيد حالة الخلية (7,7) أو اقتراح إحداثيات بديلة إذا لزم الأمر.} \\
            }
        }
    }
    & & \\ \\

    \theutterance \stepcounter{utterance}  
    & & \multicolumn{4}{p{0.6\linewidth}}{
        \cellcolor[rgb]{0.9,0.9,0.9}{
            \makecell[{{p{\linewidth}}}]{
                \texttt{\tiny{[GM$\rangle$P2]}}
                \texttt{تم تمرير رسالتك إلى اللاعب الآخر.} \\
\\ 
\\ 
\texttt{أنت تلعب حالياً الجولة 2 من أصل 12 كحد أقصى.} \\
\texttt{لقد جمعتما معاً 1 من أصل 8 عقوبة.} \\
\texttt{أرسل اللاعب الآخر هذه الرسالة:} \\
\\ 
\\ 
\texttt{"مقترح جيد. الزاوية (7,7) تبدو معقولة. سأبدأ بتحريك C من (4,4) إلى (7,7) إذا كانت الخلية فارغة. هل أنت متأكد من أن (7,7) فارغة حالياً؟ إذا كانت P لا تزال هناك، فلن أتمكن من التحرك إليها. يرجى تأكيد حالة الخلية (7,7) أو اقتراح إحداثيات بديلة إذا لزم الأمر."} \\
\\ 
\\ 
\texttt{ما هو أمرك التالي؟} \\
            }
        }
    }
    & \\ \\

    \theutterance \stepcounter{utterance}  
    & & & \multicolumn{4}{p{0.6\linewidth}}{
        \cellcolor[rgb]{0.9,0.9,0.9}{
            \makecell[{{p{\linewidth}}}]{
                \texttt{\tiny{[GM$\langle$P2]}}
                \texttt{قل: نعم، الخلية (7,7) فارغة حالياً. P موجود في (5,1) وليس في (7,7). يمكنك تحريك C إلى (7,7) بأمان. سأقوم أنا بتحريك L من (3,3) نحو (7,7) في دوراتي القادمة.} \\
            }
        }
    }
     \\ \\

    \theutterance \stepcounter{utterance}  
    & & & \multicolumn{2}{p{0.3\linewidth}}{
        \cellcolor[rgb]{0.9,0.9,0.9}{
            \makecell[{{p{\linewidth}}}]{
                \texttt{\tiny{[GM$|$GM]}}
                \texttt{قل: نعم، الخلية (7,7) فارغة حالياً. P موجود في (5,1) وليس في (7,7). يمكنك تحريك C إلى (7,7) بأمان. سأقوم أنا بتحريك L من (3,3) نحو (7,7) في دوراتي القادمة.} \\
            }
        }
    }
    & & \\ \\

    \theutterance \stepcounter{utterance}  
    & & \multicolumn{4}{p{0.6\linewidth}}{
        \cellcolor[rgb]{0.9,0.9,0.9}{
            \makecell[{{p{\linewidth}}}]{
                \texttt{\tiny{[P1$\langle$GM]}}
                \texttt{تم تمرير رسالتك إلى اللاعب الآخر.} \\
\\ 
\\ 
\texttt{أنت تلعب حالياً الجولة 2 من أصل 12 كحد أقصى.} \\
\texttt{لقد جمعتما معاً 1 من أصل 8 عقوبة.} \\
\texttt{أرسل اللاعب الآخر هذه الرسالة:} \\
\\ 
\\ 
\texttt{"نعم، الخلية (7,7) فارغة حالياً. P موجود في (5,1) وليس في (7,7). يمكنك تحريك C إلى (7,7) بأمان. سأقوم أنا بتحريك L من (3,3) نحو (7,7) في دوراتي القادمة."} \\
\\ 
\\ 
\texttt{ما هو أمرك التالي؟} \\
            }
        }
    }
    & \\ \\

    \theutterance \stepcounter{utterance}  
    & \multicolumn{4}{p{0.6\linewidth}}{
        \cellcolor[rgb]{0.9,0.9,0.9}{
            \makecell[{{p{\linewidth}}}]{
                \texttt{\tiny{[P1$\rangle$GM]}}
                \texttt{حرّوح: C, (7,7)} \\
            }
        }
    }
    & & \\ \\

    \theutterance \stepcounter{utterance}  
    & & & \multicolumn{2}{p{0.3\linewidth}}{
        \cellcolor[rgb]{0.9,0.9,0.9}{
            \makecell[{{p{\linewidth}}}]{
                \texttt{\tiny{[GM$|$GM]}}
                \texttt{حرّوح: C, (7,7)} \\
            }
        }
    }
    & & \\ \\

    \theutterance \stepcounter{utterance}  
    & & & \multicolumn{2}{p{0.3\linewidth}}{
        \cellcolor[rgb]{0.9,0.9,0.9}{
            \makecell[{{p{\linewidth}}}]{
                \texttt{\tiny{[GM$|$GM]}}
                \texttt{Invalid response format} \\
            }
        }
    }
    & & \\ \\

    \theutterance \stepcounter{utterance}  
    & & \multicolumn{4}{p{0.6\linewidth}}{
        \cellcolor[rgb]{0.9,0.9,0.9}{
            \makecell[{{p{\linewidth}}}]{
                \texttt{\tiny{[P1$\langle$GM]}}
                \texttt{عقوبة: رسالتك ليست بالتنسيق المتوقع!} \\
\\ 
\texttt{تأكد من أن ردك يحتوي فقط على `قل: <رسالة>` أو `حرّك: <كائن>, (<X>, <Y>)`، ولا شيء غير ذلك!} \\
\\ 
\texttt{لقد جمعتما معاً 2 من أصل 8 عقوبة.} \\
\texttt{يرجى المحاولة مرة أخرى!} \\
            }
        }
    }
    & \\ \\

    \theutterance \stepcounter{utterance}  
    & \multicolumn{4}{p{0.6\linewidth}}{
        \cellcolor[rgb]{0.9,0.9,0.9}{
            \makecell[{{p{\linewidth}}}]{
                \texttt{\tiny{[P1$\rangle$GM]}}
                \texttt{حرّك: C, (7,7)} \\
            }
        }
    }
    & & \\ \\

    \theutterance \stepcounter{utterance}  
    & & & \multicolumn{2}{p{0.3\linewidth}}{
        \cellcolor[rgb]{0.9,0.9,0.9}{
            \makecell[{{p{\linewidth}}}]{
                \texttt{\tiny{[GM$|$GM]}}
                \texttt{حرّك: C, (7,7)} \\
            }
        }
    }
    & & \\ \\

    \theutterance \stepcounter{utterance}  
    & & & \multicolumn{2}{p{0.3\linewidth}}{
        \cellcolor[rgb]{0.9,0.9,0.9}{
            \makecell[{{p{\linewidth}}}]{
                \texttt{\tiny{[GM$|$GM]}}
                \texttt{C moved to (7, 7)} \\
            }
        }
    }
    & & \\ \\

    \theutterance \stepcounter{utterance}  
    & & \multicolumn{4}{p{0.6\linewidth}}{
        \cellcolor[rgb]{0.9,0.9,0.9}{
            \makecell[{{p{\linewidth}}}]{
                \texttt{\tiny{[GM$\rangle$P2]}}
                \texttt{تم تمرير رسالتك إلى اللاعب الآخر.} \\
\\ 
\\ 
\texttt{أنت تلعب حالياً الجولة 3 من أصل 12 كحد أقصى.} \\
\texttt{لقد جمعتما معاً 2 من أصل 8 عقوبة.} \\
\texttt{حرّك اللاعب الآخر كائناً على شبكته.} \\
\\ 
\\ 
\texttt{ما هو أمرك التالي؟} \\
            }
        }
    }
    & \\ \\

    \theutterance \stepcounter{utterance}  
    & & & \multicolumn{4}{p{0.6\linewidth}}{
        \cellcolor[rgb]{0.9,0.9,0.9}{
            \makecell[{{p{\linewidth}}}]{
                \texttt{\tiny{[GM$\langle$P2]}}
                \texttt{حرّك: L, (7,7)} \\
            }
        }
    }
     \\ \\

    \theutterance \stepcounter{utterance}  
    & & & \multicolumn{2}{p{0.3\linewidth}}{
        \cellcolor[rgb]{0.9,0.9,0.9}{
            \makecell[{{p{\linewidth}}}]{
                \texttt{\tiny{[GM$|$GM]}}
                \texttt{حرّك: L, (7,7)} \\
            }
        }
    }
    & & \\ \\

    \theutterance \stepcounter{utterance}  
    & & & \multicolumn{2}{p{0.3\linewidth}}{
        \cellcolor[rgb]{0.9,0.9,0.9}{
            \makecell[{{p{\linewidth}}}]{
                \texttt{\tiny{[GM$|$GM]}}
                \texttt{L moved to (7, 7)} \\
            }
        }
    }
    & & \\ \\

    \theutterance \stepcounter{utterance}  
    & & \multicolumn{4}{p{0.6\linewidth}}{
        \cellcolor[rgb]{0.9,0.9,0.9}{
            \makecell[{{p{\linewidth}}}]{
                \texttt{\tiny{[P1$\langle$GM]}}
                \texttt{تم تحريك 'C' إلى (7,7) بنجاح. تبدو شبكتك المحدّثة هكذا:} \\
\texttt{```} \\
\texttt{ 1234567} \\
\texttt{╔═╤══╤══╗} \\
\texttt{║◌│◌◌│◌◌║ 1} \\
\texttt{╟─┤◌◌│◌◌║ 2} \\
\texttt{║◌│◌◌├──╢ 3} \\
\texttt{╟─┘◌◌│◌◌║ 4} \\
\texttt{║◌◌◌◌│◌◌║ 5} \\
\texttt{║◌L┌─┼──╢ 6} \\
\texttt{║◌◌│◌│PC║ 7} \\
\texttt{╚══╧═╧══╝} \\
\texttt{```} \\
\texttt{أنت تلعب حالياً الجولة 3 من أصل 12 كحد أقصى.} \\
\texttt{لقد جمعتما معاً 2 من أصل 8 عقوبة.} \\
\texttt{حرّك اللاعب الآخر كائناً على شبكته.} \\
\\ 
\\ 
\texttt{ما هو أمرك التالي؟} \\
            }
        }
    }
    & \\ \\

    \theutterance \stepcounter{utterance}  
    & \multicolumn{4}{p{0.6\linewidth}}{
        \cellcolor[rgb]{0.9,0.9,0.9}{
            \makecell[{{p{\linewidth}}}]{
                \texttt{\tiny{[P1$\rangle$GM]}}
                \texttt{قل: ممتاز! لقد تحرك C بنجاح إلى (7,7). الآن أرى L في (6,6) و P في (7,7)؟ انتظر، اللوحة تظهر P و C في نفس الصف 7. دعني أتحقق: C في (7,7). P في (7,7) أيضاً؟ لا، اللوحة تظهر `PC` في الخلية (7,7)؟ لا، هذا غير ممكن. دعني أقرأ اللوحة بعناية:} \\
\texttt{الصف 7: `║◌◌│◌│PC║`} \\
\texttt{هذا يعني أن العمود 6 يحتوي على P؟ لا، الحدود واضحة.} \\
\texttt{الخلايا في الصف 7:} \\
\texttt{العمود 1: ◌} \\
\texttt{العمود 2: ◌} \\
\texttt{العمود 3: ◌ (بعد الخط العمودي)} \\
\texttt{العمود 4: ◌} \\
\texttt{العمود 5: ◌} \\
\texttt{العمود 6: P} \\
\texttt{العمود 7: C} \\
\\ 
\texttt{انتظر، الرمز `PC` يظهر في العمود 6 و 7؟} \\
\texttt{لننظر إلى الهيكل:} \\
\texttt{`║◌◌│◌│PC║`} \\
\texttt{العمود 1: ◌} \\
\texttt{العمود 2: ◌} \\
\texttt{الفاصل} \\
\texttt{العمود 3: ◌} \\
\texttt{الفاصل} \\
\texttt{العمود 4: ◌} \\
\texttt{الفاصل} \\
\texttt{العمود 5: ◌} \\
\texttt{الفاصل} \\
\texttt{العمود 6: P} \\
\texttt{العمود 7: C} \\
\\ 
\texttt{إذاً، P في (6,7) و C في (7,7).} \\
\texttt{L في (6,6).} \\
\\ 
\texttt{الهدف هو تجميعهم في (7,7).} \\
\texttt{C موجود بالفعل في (7,7).} \\
\texttt{L في (6,6).} \\
\texttt{P في (6,7).} \\
\\ 
\texttt{يجب تحريك L و P إلى (7,7). لكن (7,7) مشغولة بـ C.} \\
\texttt{هل يجب أن نحرك C أولاً؟ لا، الهدف هو التجميع.} \\
\texttt{ربما يجب أن نحرك C إلى مكان آخر مؤقتاً؟ أو ربما الهدف ليس (7,7) بالضبط بل منطقة قريبة؟} \\
\texttt{اللاعب الآخر قال إن (7,7) فارغة، لكن C تحرك إليها.} \\
\texttt{الآن C في (7,7).} \\
\texttt{L في (6,6).} \\
\texttt{P في (6,7).} \\
\\ 
\texttt{لنحاول تحريك L إلى (7,6) ثم إلى (7,7)؟ لا، (7,7) مشغولة.} \\
\texttt{ربما نحرك C إلى (7,6)؟} \\
\texttt{أو نحرك P إلى (7,7)؟ لا، مشغولة.} \\
\\ 
\texttt{دعنا نغير الاستراتيجية قليلاً. لنحاول تجميعهم في (6,6)؟} \\
\texttt{L في (6,6).} \\
\texttt{P في (6,7).} \\
\texttt{C في (7,7).} \\
\\ 
\texttt{إذا حركنا P إلى (6,6)، لا يمكن لأن L هناك.} \\
\texttt{إذا حركنا L إلى (6,7)، لا يمكن لأن P هناك.} \\
\\ 
\texttt{ربما نحرك C إلى (7,6)؟} \\
\texttt{ثم نحرك P إلى (7,7)؟} \\
\texttt{ثم نحرك L إلى (7,6)؟ لا، C هناك.} \\
\\ 
\texttt{دعنا نحاول تحريك L إلى (7,6).} \\
\texttt{الخلية (7,6) فارغة حالياً؟} \\
\texttt{الصف 7: العمود 6 هو P؟ لا، P في (6,7)؟} \\
\texttt{لنعد قراءة اللوحة:} \\
\texttt{`║◌◌│◌│PC║`} \\
\texttt{الخلايا:} \\
\texttt{(7,1): ◌} \\
\texttt{(7,2): ◌} \\
\texttt{(7,3): ◌} \\
\texttt{(7,4): ◌} \\
\texttt{(7,5): ◌} \\
\texttt{(7,6): P} \\
\texttt{(7,7): C} \\
\\ 
\texttt{إذاً P في (7,6) و C في (7,7).} \\
\texttt{L في (6,6).} \\
\\ 
\texttt{الآن:} \\
\texttt{L في (6,6).} \\
\texttt{P في (7,6).} \\
\texttt{C في (7,7).} \\
\\ 
\texttt{نريد تجميعهم.} \\
\texttt{يمكننا تحريك L إلى (7,6)؟ لا، P هناك.} \\
\texttt{يمكننا تحريك P إلى (7,7)؟ لا، C هناك.} \\
\texttt{يمكننا تحريك C إلى (7,6)؟ لا، P هناك.} \\
\\ 
\texttt{يجب تحريك واحد منهم بعيداً.} \\
\texttt{لنحرك C إلى (7,5)؟} \\
\texttt{أو لنحرك P إلى (6,6)؟ لا، L هناك.} \\
\\ 
\texttt{لنحرك L إلى (6,7)؟ لا، P في (7,6)؟ لا، P في (7,6).} \\
\texttt{الخلية (6,7) فارغة؟} \\
\texttt{الصف 6: `║◌L┌─┼──╢`} \\
\texttt{(6,1): ◌} \\
\texttt{(6,2): L} \\
\texttt{(6,3): ◌} \\
\texttt{(6,4): ◌} \\
\texttt{(6,5): ◌} \\
\texttt{(6,6): ◌ (لأن L في (6,2)؟ لا، L في (6,2)؟} \\
\texttt{لنقرأ الصف 6: `║◌L┌─┼──╢`} \\
\texttt{العمود 1: ◌} \\
\texttt{العمود 2: L} \\
\texttt{العمود 3: ◌} \\
\texttt{العمود 4: ◌} \\
\texttt{العمود 5: ◌} \\
\texttt{العمود 6: ◌} \\
\texttt{العمود 7: ◌} \\
\\ 
\texttt{إذاً L في (2,6)؟ لا، الإحداثيات X هي الأعمدة، Y هي الصفوف.} \\
\texttt{الصف 6، العمود 2 هو L. إذن L في (2,6).} \\
\texttt{الصف 7، العمود 6 هو P. إذن P في (6,7).} \\
\texttt{الصف 7، العمود 7 هو C. إذن C في (7,7).} \\
\\ 
\texttt{الآن:} \\
\texttt{L في (2,6).} \\
\texttt{P في (6,7).} \\
\texttt{C في (7,7).} \\
\\ 
\texttt{نريد تجميعهم في (7,7).} \\
\texttt{C موجود في (7,7).} \\
\texttt{P في (6,7). يمكن تحريك P إلى (7,7)؟ لا، C هناك.} \\
\texttt{L في (2,6). بعيد.} \\
\\ 
\texttt{لنحرك C إلى (7,6)؟} \\
\texttt{الخلية (7,6) مشغولة بـ P.} \\
\texttt{لنحرك P إلى (7,7)؟ لا، C هناك.} \\
\\ 
\texttt{لنحرك C إلى (6,7)؟ لا، P في (6,7).} \\
\\ 
\texttt{يجب تحريك C أو P.} \\
\texttt{لنحرك C إلى (7,5)؟} \\
\texttt{الخلية (7,5) فارغة؟} \\
\texttt{الصف 7: `║◌◌│◌│PC║`} \\
\texttt{(7,1): ◌} \\
\texttt{(7,2): ◌} \\
\texttt{(7,3): ◌} \\
\texttt{(7,4): ◌} \\
\texttt{(7,5): ◌} \\
\texttt{(7,6): P} \\
\texttt{(7,7): C} \\
\\ 
\texttt{نعم، (7,5) فارغة.} \\
\texttt{لنحرك C إلى (7,5).} \\
\texttt{ثم نحرك P إلى (7,7).} \\
\texttt{ثم نحرك L إلى (7,7)؟ لا، P هناك.} \\
\texttt{ثم نحرك P إلى (7,6)؟} \\
\texttt{ثم نحرك L إلى (7,7).} \\
\\ 
\texttt{هذا معقد.} \\
\texttt{ربما نجمعهم في (6,6)؟} \\
\texttt{L في (2,6).} \\
\texttt{P في (6,7).} \\
\texttt{C في (7,7).} \\
\\ 
\texttt{لنحرك C إلى (6,7)؟ لا، P هناك.} \\
\texttt{لنحرك P إلى (6,6)؟} \\
\texttt{الخلية (6,6) فارغة؟} \\
\texttt{الصف 6: `║◌L┌─┼──╢`} \\
\texttt{(6,1): ◌} \\
\texttt{(6,2): L} \\
\texttt{(6,3): ◌} \\
\texttt{(6,4): ◌} \\
\texttt{(6,5): ◌} \\
\texttt{(6,6): ◌} \\
\texttt{(6,7): ◌} \\
\\ 
\texttt{نعم، (6,6) فارغة.} \\
\texttt{لنحرك P من (6,7) إلى (6,6).} \\
\texttt{ثم نحرك C من (7,7) إلى (6,7).} \\
\texttt{ثم نحرك L من (2,6) إلى (6,6)؟ لا، P هناك.} \\
\texttt{ثم نحرك P إلى (6,7)؟ لا، C هناك.} \\
\\ 
\texttt{دعنا نحاول تجميعهم في (7,7) مرة أخرى.} \\
\texttt{C في (7,7).} \\
\texttt{P في (6,7).} \\
\texttt{L في (2,6).} \\
\\ 
\texttt{لنحرك P إلى (7,7)؟ لا، C هناك.} \\
\texttt{لنحرك C إلى (7,6)؟ لا، P هناك.} \\
\\ 
\texttt{لنحرك C إلى (7,5).} \\
\texttt{حرّك: C, (7,5)} \\
            }
        }
    }
    & & \\ \\

    \theutterance \stepcounter{utterance}  
    & & & \multicolumn{2}{p{0.3\linewidth}}{
        \cellcolor[rgb]{0.9,0.9,0.9}{
            \makecell[{{p{\linewidth}}}]{
                \texttt{\tiny{[GM$|$GM]}}
                \texttt{قل: ممتاز! لقد تحرك C بنجاح إلى (7,7). الآن أرى L في (6,6) و P في (7,7)؟ انتظر، اللوحة تظهر P و C في نفس الصف 7. دعني أتحقق: C في (7,7). P في (7,7) أيضاً؟ لا، اللوحة تظهر `PC` في الخلية (7,7)؟ لا، هذا غير ممكن. دعني أقرأ اللوحة بعناية:} \\
\texttt{الصف 7: `║◌◌│◌│PC║`} \\
\texttt{هذا يعني أن العمود 6 يحتوي على P؟ لا، الحدود واضحة.} \\
\texttt{الخلايا في الصف 7:} \\
\texttt{العمود 1: ◌} \\
\texttt{العمود 2: ◌} \\
\texttt{العمود 3: ◌ (بعد الخط العمودي)} \\
\texttt{العمود 4: ◌} \\
\texttt{العمود 5: ◌} \\
\texttt{العمود 6: P} \\
\texttt{العمود 7: C} \\
\\ 
\texttt{انتظر، الرمز `PC` يظهر في العمود 6 و 7؟} \\
\texttt{لننظر إلى الهيكل:} \\
\texttt{`║◌◌│◌│PC║`} \\
\texttt{العمود 1: ◌} \\
\texttt{العمود 2: ◌} \\
\texttt{الفاصل} \\
\texttt{العمود 3: ◌} \\
\texttt{الفاصل} \\
\texttt{العمود 4: ◌} \\
\texttt{الفاصل} \\
\texttt{العمود 5: ◌} \\
\texttt{الفاصل} \\
\texttt{العمود 6: P} \\
\texttt{العمود 7: C} \\
\\ 
\texttt{إذاً، P في (6,7) و C في (7,7).} \\
\texttt{L في (6,6).} \\
\\ 
\texttt{الهدف هو تجميعهم في (7,7).} \\
\texttt{C موجود بالفعل في (7,7).} \\
\texttt{L في (6,6).} \\
\texttt{P في (6,7).} \\
\\ 
\texttt{يجب تحريك L و P إلى (7,7). لكن (7,7) مشغولة بـ C.} \\
\texttt{هل يجب أن نحرك C أولاً؟ لا، الهدف هو التجميع.} \\
\texttt{ربما يجب أن نحرك C إلى مكان آخر مؤقتاً؟ أو ربما الهدف ليس (7,7) بالضبط بل منطقة قريبة؟} \\
\texttt{اللاعب الآخر قال إن (7,7) فارغة، لكن C تحرك إليها.} \\
\texttt{الآن C في (7,7).} \\
\texttt{L في (6,6).} \\
\texttt{P في (6,7).} \\
\\ 
\texttt{لنحاول تحريك L إلى (7,6) ثم إلى (7,7)؟ لا، (7,7) مشغولة.} \\
\texttt{ربما نحرك C إلى (7,6)؟} \\
\texttt{أو نحرك P إلى (7,7)؟ لا، مشغولة.} \\
\\ 
\texttt{دعنا نغير الاستراتيجية قليلاً. لنحاول تجميعهم في (6,6)؟} \\
\texttt{L في (6,6).} \\
\texttt{P في (6,7).} \\
\texttt{C في (7,7).} \\
\\ 
\texttt{إذا حركنا P إلى (6,6)، لا يمكن لأن L هناك.} \\
\texttt{إذا حركنا L إلى (6,7)، لا يمكن لأن P هناك.} \\
\\ 
\texttt{ربما نحرك C إلى (7,6)؟} \\
\texttt{ثم نحرك P إلى (7,7)؟} \\
\texttt{ثم نحرك L إلى (7,6)؟ لا، C هناك.} \\
\\ 
\texttt{دعنا نحاول تحريك L إلى (7,6).} \\
\texttt{الخلية (7,6) فارغة حالياً؟} \\
\texttt{الصف 7: العمود 6 هو P؟ لا، P في (6,7)؟} \\
\texttt{لنعد قراءة اللوحة:} \\
\texttt{`║◌◌│◌│PC║`} \\
\texttt{الخلايا:} \\
\texttt{(7,1): ◌} \\
\texttt{(7,2): ◌} \\
\texttt{(7,3): ◌} \\
\texttt{(7,4): ◌} \\
\texttt{(7,5): ◌} \\
\texttt{(7,6): P} \\
\texttt{(7,7): C} \\
\\ 
\texttt{إذاً P في (7,6) و C في (7,7).} \\
\texttt{L في (6,6).} \\
\\ 
\texttt{الآن:} \\
\texttt{L في (6,6).} \\
\texttt{P في (7,6).} \\
\texttt{C في (7,7).} \\
\\ 
\texttt{نريد تجميعهم.} \\
\texttt{يمكننا تحريك L إلى (7,6)؟ لا، P هناك.} \\
\texttt{يمكننا تحريك P إلى (7,7)؟ لا، C هناك.} \\
\texttt{يمكننا تحريك C إلى (7,6)؟ لا، P هناك.} \\
\\ 
\texttt{يجب تحريك واحد منهم بعيداً.} \\
\texttt{لنحرك C إلى (7,5)؟} \\
\texttt{أو لنحرك P إلى (6,6)؟ لا، L هناك.} \\
\\ 
\texttt{لنحرك L إلى (6,7)؟ لا، P في (7,6)؟ لا، P في (7,6).} \\
\texttt{الخلية (6,7) فارغة؟} \\
\texttt{الصف 6: `║◌L┌─┼──╢`} \\
\texttt{(6,1): ◌} \\
\texttt{(6,2): L} \\
\texttt{(6,3): ◌} \\
\texttt{(6,4): ◌} \\
\texttt{(6,5): ◌} \\
\texttt{(6,6): ◌ (لأن L في (6,2)؟ لا، L في (6,2)؟} \\
\texttt{لنقرأ الصف 6: `║◌L┌─┼──╢`} \\
\texttt{العمود 1: ◌} \\
\texttt{العمود 2: L} \\
\texttt{العمود 3: ◌} \\
\texttt{العمود 4: ◌} \\
\texttt{العمود 5: ◌} \\
\texttt{العمود 6: ◌} \\
\texttt{العمود 7: ◌} \\
\\ 
\texttt{إذاً L في (2,6)؟ لا، الإحداثيات X هي الأعمدة، Y هي الصفوف.} \\
\texttt{الصف 6، العمود 2 هو L. إذن L في (2,6).} \\
\texttt{الصف 7، العمود 6 هو P. إذن P في (6,7).} \\
\texttt{الصف 7، العمود 7 هو C. إذن C في (7,7).} \\
\\ 
\texttt{الآن:} \\
\texttt{L في (2,6).} \\
\texttt{P في (6,7).} \\
\texttt{C في (7,7).} \\
\\ 
\texttt{نريد تجميعهم في (7,7).} \\
\texttt{C موجود في (7,7).} \\
\texttt{P في (6,7). يمكن تحريك P إلى (7,7)؟ لا، C هناك.} \\
\texttt{L في (2,6). بعيد.} \\
\\ 
\texttt{لنحرك C إلى (7,6)؟} \\
\texttt{الخلية (7,6) مشغولة بـ P.} \\
\texttt{لنحرك P إلى (7,7)؟ لا، C هناك.} \\
\\ 
\texttt{لنحرك C إلى (6,7)؟ لا، P في (6,7).} \\
\\ 
\texttt{يجب تحريك C أو P.} \\
\texttt{لنحرك C إلى (7,5)؟} \\
\texttt{الخلية (7,5) فارغة؟} \\
\texttt{الصف 7: `║◌◌│◌│PC║`} \\
\texttt{(7,1): ◌} \\
\texttt{(7,2): ◌} \\
\texttt{(7,3): ◌} \\
\texttt{(7,4): ◌} \\
\texttt{(7,5): ◌} \\
\texttt{(7,6): P} \\
\texttt{(7,7): C} \\
\\ 
\texttt{نعم، (7,5) فارغة.} \\
\texttt{لنحرك C إلى (7,5).} \\
\texttt{ثم نحرك P إلى (7,7).} \\
\texttt{ثم نحرك L إلى (7,7)؟ لا، P هناك.} \\
\texttt{ثم نحرك P إلى (7,6)؟} \\
\texttt{ثم نحرك L إلى (7,7).} \\
\\ 
\texttt{هذا معقد.} \\
\texttt{ربما نجمعهم في (6,6)؟} \\
\texttt{L في (2,6).} \\
\texttt{P في (6,7).} \\
\texttt{C في (7,7).} \\
\\ 
\texttt{لنحرك C إلى (6,7)؟ لا، P هناك.} \\
\texttt{لنحرك P إلى (6,6)؟} \\
\texttt{الخلية (6,6) فارغة؟} \\
\texttt{الصف 6: `║◌L┌─┼──╢`} \\
\texttt{(6,1): ◌} \\
\texttt{(6,2): L} \\
\texttt{(6,3): ◌} \\
\texttt{(6,4): ◌} \\
\texttt{(6,5): ◌} \\
\texttt{(6,6): ◌} \\
\texttt{(6,7): ◌} \\
\\ 
\texttt{نعم، (6,6) فارغة.} \\
\texttt{لنحرك P من (6,7) إلى (6,6).} \\
\texttt{ثم نحرك C من (7,7) إلى (6,7).} \\
\texttt{ثم نحرك L من (2,6) إلى (6,6)؟ لا، P هناك.} \\
\texttt{ثم نحرك P إلى (6,7)؟ لا، C هناك.} \\
\\ 
\texttt{دعنا نحاول تجميعهم في (7,7) مرة أخرى.} \\
\texttt{C في (7,7).} \\
\texttt{P في (6,7).} \\
\texttt{L في (2,6).} \\
\\ 
\texttt{لنحرك P إلى (7,7)؟ لا، C هناك.} \\
\texttt{لنحرك C إلى (7,6)؟ لا، P هناك.} \\
\\ 
\texttt{لنحرك C إلى (7,5).} \\
\texttt{حرّك: C, (7,5)} \\
            }
        }
    }
    & & \\ \\

    \theutterance \stepcounter{utterance}  
    & & & \multicolumn{2}{p{0.3\linewidth}}{
        \cellcolor[rgb]{0.9,0.9,0.9}{
            \makecell[{{p{\linewidth}}}]{
                \texttt{\tiny{[GM$|$GM]}}
                \texttt{Invalid response format: قل: ممتاز! لقد تحرك C بنجاح إلى (7,7). الآن أرى L في (6,6) و P في (7,7)؟ انتظر، اللوحة تظهر P و C في نفس الصف 7. دعني أتحقق: C في (7,7). P في (7,7) أيضاً؟ لا، اللوحة تظهر PC في الخلية (7,7)؟ لا، هذا غير ممكن. دعني أقرأ اللوحة بعناية:} \\
\texttt{الصف 7: ║◌◌│◌│PC║} \\
\texttt{هذا يعني أن العمود 6 يحتوي على P؟ لا، الحدود واضحة.} \\
\texttt{الخلايا في الصف 7:} \\
\texttt{العمود 1: ◌} \\
\texttt{العمود 2: ◌} \\
\texttt{العمود 3: ◌ (بعد الخط العمودي)} \\
\texttt{العمود 4: ◌} \\
\texttt{العمود 5: ◌} \\
\texttt{العمود 6: P} \\
\texttt{العمود 7: C} \\
\\ 
\texttt{انتظر، الرمز PC يظهر في العمود 6 و 7؟} \\
\texttt{لننظر إلى الهيكل:} \\
\texttt{║◌◌│◌│PC║} \\
\texttt{العمود 1: ◌} \\
\texttt{العمود 2: ◌} \\
\texttt{الفاصل} \\
\texttt{العمود 3: ◌} \\
\texttt{الفاصل} \\
\texttt{العمود 4: ◌} \\
\texttt{الفاصل} \\
\texttt{العمود 5: ◌} \\
\texttt{الفاصل} \\
\texttt{العمود 6: P} \\
\texttt{العمود 7: C} \\
\\ 
\texttt{إذاً، P في (6,7) و C في (7,7).} \\
\texttt{L في (6,6).} \\
\\ 
\texttt{الهدف هو تجميعهم في (7,7).} \\
\texttt{C موجود بالفعل في (7,7).} \\
\texttt{L في (6,6).} \\
\texttt{P في (6,7).} \\
\\ 
\texttt{يجب تحريك L و P إلى (7,7). لكن (7,7) مشغولة بـ C.} \\
\texttt{هل يجب أن نحرك C أولاً؟ لا، الهدف هو التجميع.} \\
\texttt{ربما يجب أن نحرك C إلى مكان آخر مؤقتاً؟ أو ربما الهدف ليس (7,7) بالضبط بل منطقة قريبة؟} \\
\texttt{اللاعب الآخر قال إن (7,7) فارغة، لكن C تحرك إليها.} \\
\texttt{الآن C في (7,7).} \\
\texttt{L في (6,6).} \\
\texttt{P في (6,7).} \\
\\ 
\texttt{لنحاول تحريك L إلى (7,6) ثم إلى (7,7)؟ لا، (7,7) مشغولة.} \\
\texttt{ربما نحرك C إلى (7,6)؟} \\
\texttt{أو نحرك P إلى (7,7)؟ لا، مشغولة.} \\
\\ 
\texttt{دعنا نغير الاستراتيجية قليلاً. لنحاول تجميعهم في (6,6)؟} \\
\texttt{L في (6,6).} \\
\texttt{P في (6,7).} \\
\texttt{C في (7,7).} \\
\\ 
\texttt{إذا حركنا P إلى (6,6)، لا يمكن لأن L هناك.} \\
\texttt{إذا حركنا L إلى (6,7)، لا يمكن لأن P هناك.} \\
\\ 
\texttt{ربما نحرك C إلى (7,6)؟} \\
\texttt{ثم نحرك P إلى (7,7)؟} \\
\texttt{ثم نحرك L إلى (7,6)؟ لا، C هناك.} \\
\\ 
\texttt{دعنا نحاول تحريك L إلى (7,6).} \\
\texttt{الخلية (7,6) فارغة حالياً؟} \\
\texttt{الصف 7: العمود 6 هو P؟ لا، P في (6,7)؟} \\
\texttt{لنعد قراءة اللوحة:} \\
\texttt{║◌◌│◌│PC║} \\
\texttt{الخلايا:} \\
\texttt{(7,1): ◌} \\
\texttt{(7,2): ◌} \\
\texttt{(7,3): ◌} \\
\texttt{(7,4): ◌} \\
\texttt{(7,5): ◌} \\
\texttt{(7,6): P} \\
\texttt{(7,7): C} \\
\\ 
\texttt{إذاً P في (7,6) و C في (7,7).} \\
\texttt{L في (6,6).} \\
\\ 
\texttt{الآن:} \\
\texttt{L في (6,6).} \\
\texttt{P في (7,6).} \\
\texttt{C في (7,7).} \\
\\ 
\texttt{نريد تجميعهم.} \\
\texttt{يمكننا تحريك L إلى (7,6)؟ لا، P هناك.} \\
\texttt{يمكننا تحريك P إلى (7,7)؟ لا، C هناك.} \\
\texttt{يمكننا تحريك C إلى (7,6)؟ لا، P هناك.} \\
\\ 
\texttt{يجب تحريك واحد منهم بعيداً.} \\
\texttt{لنحرك C إلى (7,5)؟} \\
\texttt{أو لنحرك P إلى (6,6)؟ لا، L هناك.} \\
\\ 
\texttt{لنحرك L إلى (6,7)؟ لا، P في (7,6)؟ لا، P في (7,6).} \\
\texttt{الخلية (6,7) فارغة؟} \\
\texttt{الصف 6: ║◌L┌─┼──╢} \\
\texttt{(6,1): ◌} \\
\texttt{(6,2): L} \\
\texttt{(6,3): ◌} \\
\texttt{(6,4): ◌} \\
\texttt{(6,5): ◌} \\
\texttt{(6,6): ◌ (لأن L في (6,2)؟ لا، L في (6,2)؟} \\
\texttt{لنقرأ الصف 6: ║◌L┌─┼──╢} \\
\texttt{العمود 1: ◌} \\
\texttt{العمود 2: L} \\
\texttt{العمود 3: ◌} \\
\texttt{العمود 4: ◌} \\
\texttt{العمود 5: ◌} \\
\texttt{العمود 6: ◌} \\
\texttt{العمود 7: ◌} \\
\\ 
\texttt{إذاً L في (2,6)؟ لا، الإحداثيات X هي الأعمدة، Y هي الصفوف.} \\
\texttt{الصف 6، العمود 2 هو L. إذن L في (2,6).} \\
\texttt{الصف 7، العمود 6 هو P. إذن P في (6,7).} \\
\texttt{الصف 7، العمود 7 هو C. إذن C في (7,7).} \\
\\ 
\texttt{الآن:} \\
\texttt{L في (2,6).} \\
\texttt{P في (6,7).} \\
\texttt{C في (7,7).} \\
\\ 
\texttt{نريد تجميعهم في (7,7).} \\
\texttt{C موجود في (7,7).} \\
\texttt{P في (6,7). يمكن تحريك P إلى (7,7)؟ لا، C هناك.} \\
\texttt{L في (2,6). بعيد.} \\
\\ 
\texttt{لنحرك C إلى (7,6)؟} \\
\texttt{الخلية (7,6) مشغولة بـ P.} \\
\texttt{لنحرك P إلى (7,7)؟ لا، C هناك.} \\
\\ 
\texttt{لنحرك C إلى (6,7)؟ لا، P في (6,7).} \\
\\ 
\texttt{يجب تحريك C أو P.} \\
\texttt{لنحرك C إلى (7,5)؟} \\
\texttt{الخلية (7,5) فارغة؟} \\
\texttt{الصف 7: ║◌◌│◌│PC║} \\
\texttt{(7,1): ◌} \\
\texttt{(7,2): ◌} \\
\texttt{(7,3): ◌} \\
\texttt{(7,4): ◌} \\
\texttt{(7,5): ◌} \\
\texttt{(7,6): P} \\
\texttt{(7,7): C} \\
\\ 
\texttt{نعم، (7,5) فارغة.} \\
\texttt{لنحرك C إلى (7,5).} \\
\texttt{ثم نحرك P إلى (7,7).} \\
\texttt{ثم نحرك L إلى (7,7)؟ لا، P هناك.} \\
\texttt{ثم نحرك P إلى (7,6)؟} \\
\texttt{ثم نحرك L إلى (7,7).} \\
\\ 
\texttt{هذا معقد.} \\
\texttt{ربما نجمعهم في (6,6)؟} \\
\texttt{L في (2,6).} \\
\texttt{P في (6,7).} \\
\texttt{C في (7,7).} \\
\\ 
\texttt{لنحرك C إلى (6,7)؟ لا، P هناك.} \\
\texttt{لنحرك P إلى (6,6)؟} \\
\texttt{الخلية (6,6) فارغة؟} \\
\texttt{الصف 6: ║◌L┌─┼──╢} \\
\texttt{(6,1): ◌} \\
\texttt{(6,2): L} \\
\texttt{(6,3): ◌} \\
\texttt{(6,4): ◌} \\
\texttt{(6,5): ◌} \\
\texttt{(6,6): ◌} \\
\texttt{(6,7): ◌} \\
\\ 
\texttt{نعم، (6,6) فارغة.} \\
\texttt{لنحرك P من (6,7) إلى (6,6).} \\
\texttt{ثم نحرك C من (7,7) إلى (6,7).} \\
\texttt{ثم نحرك L من (2,6) إلى (6,6)؟ لا، P هناك.} \\
\texttt{ثم نحرك P إلى (6,7)؟ لا، C هناك.} \\
\\ 
\texttt{دعنا نحاول تجميعهم في (7,7) مرة أخرى.} \\
\texttt{C في (7,7).} \\
\texttt{P في (6,7).} \\
\texttt{L في (2,6).} \\
\\ 
\texttt{لنحرك P إلى (7,7)؟ لا، C هناك.} \\
\texttt{لنحرك C إلى (7,6)؟ لا، P هناك.} \\
\\ 
\texttt{لنحرك C إلى (7,5).} \\
\texttt{حرّك: C, (7,5)} \\
            }
        }
    }
    & & \\ \\

    \theutterance \stepcounter{utterance}  
    & & \multicolumn{4}{p{0.6\linewidth}}{
        \cellcolor[rgb]{0.9,0.9,0.9}{
            \makecell[{{p{\linewidth}}}]{
                \texttt{\tiny{[P1$\langle$GM]}}
                \texttt{عقوبة: تحتوي رسالتك على أكثر من أمر واحد!} \\
\\ 
\texttt{تأكد من أن ردك يحتوي فقط على `قل: <رسالة>` أو `حرّك: <كائن>, (<X>, <Y>)`، ولا شيء غير ذلك!} \\
\\ 
\texttt{لقد جمعتما معاً 3 من أصل 8 عقوبة.} \\
\texttt{يرجى المحاولة مرة أخرى!} \\
            }
        }
    }
    & \\ \\

    \theutterance \stepcounter{utterance}  
    & \multicolumn{4}{p{0.6\linewidth}}{
        \cellcolor[rgb]{0.9,0.9,0.9}{
            \makecell[{{p{\linewidth}}}]{
                \texttt{\tiny{[P1$\rangle$GM]}}
                \texttt{حرّك: C, (7,5)} \\
            }
        }
    }
    & & \\ \\

    \theutterance \stepcounter{utterance}  
    & & & \multicolumn{2}{p{0.3\linewidth}}{
        \cellcolor[rgb]{0.9,0.9,0.9}{
            \makecell[{{p{\linewidth}}}]{
                \texttt{\tiny{[GM$|$GM]}}
                \texttt{حرّك: C, (7,5)} \\
            }
        }
    }
    & & \\ \\

    \theutterance \stepcounter{utterance}  
    & & & \multicolumn{2}{p{0.3\linewidth}}{
        \cellcolor[rgb]{0.9,0.9,0.9}{
            \makecell[{{p{\linewidth}}}]{
                \texttt{\tiny{[GM$|$GM]}}
                \texttt{C moved to (7, 5)} \\
            }
        }
    }
    & & \\ \\

    \theutterance \stepcounter{utterance}  
    & & \multicolumn{4}{p{0.6\linewidth}}{
        \cellcolor[rgb]{0.9,0.9,0.9}{
            \makecell[{{p{\linewidth}}}]{
                \texttt{\tiny{[GM$\rangle$P2]}}
                \texttt{تم تحريك 'L' إلى (7,7) بنجاح. تبدو شبكتك المحدّثة هكذا:} \\
\texttt{```} \\
\texttt{ 1234567} \\
\texttt{╔═╤══╤══╗} \\
\texttt{║◌│◌◌│◌◌║ 1} \\
\texttt{╟─┤◌◌│◌◌║ 2} \\
\texttt{║◌│◌◌├──╢ 3} \\
\texttt{╟─┘◌◌│◌C║ 4} \\
\texttt{║P◌◌◌│◌◌║ 5} \\
\texttt{║◌◌┌─┼──╢ 6} \\
\texttt{║◌◌│◌│◌L║ 7} \\
\texttt{╚══╧═╧══╝} \\
\texttt{```} \\
\texttt{أنت تلعب حالياً الجولة 4 من أصل 12 كحد أقصى.} \\
\texttt{لقد جمعتما معاً 3 من أصل 8 عقوبة.} \\
\texttt{حرّك اللاعب الآخر كائناً على شبكته.} \\
\\ 
\\ 
\texttt{ما هو أمرك التالي؟} \\
            }
        }
    }
    & \\ \\

    \theutterance \stepcounter{utterance}  
    & & & \multicolumn{4}{p{0.6\linewidth}}{
        \cellcolor[rgb]{0.9,0.9,0.9}{
            \makecell[{{p{\linewidth}}}]{
                \texttt{\tiny{[GM$\langle$P2]}}
                \texttt{حرّح: C, (7,7)} \\
            }
        }
    }
     \\ \\

    \theutterance \stepcounter{utterance}  
    & & & \multicolumn{2}{p{0.3\linewidth}}{
        \cellcolor[rgb]{0.9,0.9,0.9}{
            \makecell[{{p{\linewidth}}}]{
                \texttt{\tiny{[GM$|$GM]}}
                \texttt{حرّح: C, (7,7)} \\
            }
        }
    }
    & & \\ \\

    \theutterance \stepcounter{utterance}  
    & & & \multicolumn{2}{p{0.3\linewidth}}{
        \cellcolor[rgb]{0.9,0.9,0.9}{
            \makecell[{{p{\linewidth}}}]{
                \texttt{\tiny{[GM$|$GM]}}
                \texttt{Invalid response format} \\
            }
        }
    }
    & & \\ \\

    \theutterance \stepcounter{utterance}  
    & & \multicolumn{4}{p{0.6\linewidth}}{
        \cellcolor[rgb]{0.9,0.9,0.9}{
            \makecell[{{p{\linewidth}}}]{
                \texttt{\tiny{[GM$\rangle$P2]}}
                \texttt{عقوبة: رسالتك ليست بالتنسيق المتوقع!} \\
\\ 
\texttt{تأكد من أن ردك يحتوي فقط على `قل: <رسالة>` أو `حرّك: <كائن>, (<X>, <Y>)`، ولا شيء غير ذلك!} \\
\\ 
\texttt{لقد جمعتما معاً 4 من أصل 8 عقوبة.} \\
\texttt{يرجى المحاولة مرة أخرى!} \\
            }
        }
    }
    & \\ \\

    \theutterance \stepcounter{utterance}  
    & & & \multicolumn{4}{p{0.6\linewidth}}{
        \cellcolor[rgb]{0.9,0.9,0.9}{
            \makecell[{{p{\linewidth}}}]{
                \texttt{\tiny{[GM$\langle$P2]}}
                \texttt{حرّك: C, (7,7)} \\
            }
        }
    }
     \\ \\

    \theutterance \stepcounter{utterance}  
    & & & \multicolumn{2}{p{0.3\linewidth}}{
        \cellcolor[rgb]{0.9,0.9,0.9}{
            \makecell[{{p{\linewidth}}}]{
                \texttt{\tiny{[GM$|$GM]}}
                \texttt{حرّك: C, (7,7)} \\
            }
        }
    }
    & & \\ \\

    \theutterance \stepcounter{utterance}  
    & & & \multicolumn{2}{p{0.3\linewidth}}{
        \cellcolor[rgb]{0.9,0.9,0.9}{
            \makecell[{{p{\linewidth}}}]{
                \texttt{\tiny{[GM$|$GM]}}
                \texttt{عقوبة: (7,7) ليست فارغة، بل تحتوي على 'L'.} \\
\\ 
\texttt{لقد جمعتما معاً 5 من أصل 8 عقوبة.} \\
\texttt{يرجى المحاولة مرة أخرى!} \\
            }
        }
    }
    & & \\ \\

    \theutterance \stepcounter{utterance}  
    & & \multicolumn{4}{p{0.6\linewidth}}{
        \cellcolor[rgb]{0.9,0.9,0.9}{
            \makecell[{{p{\linewidth}}}]{
                \texttt{\tiny{[GM$\rangle$P2]}}
                \texttt{عقوبة: (7,7) ليست فارغة، بل تحتوي على 'L'.} \\
\\ 
\texttt{لقد جمعتما معاً 5 من أصل 8 عقوبة.} \\
\texttt{يرجى المحاولة مرة أخرى!} \\
            }
        }
    }
    & \\ \\

    \theutterance \stepcounter{utterance}  
    & & & \multicolumn{4}{p{0.6\linewidth}}{
        \cellcolor[rgb]{0.9,0.9,0.9}{
            \makecell[{{p{\linewidth}}}]{
                \texttt{\tiny{[GM$\langle$P2]}}
                \texttt{حرّك: C, (4,7)} \\
            }
        }
    }
     \\ \\

    \theutterance \stepcounter{utterance}  
    & & & \multicolumn{2}{p{0.3\linewidth}}{
        \cellcolor[rgb]{0.9,0.9,0.9}{
            \makecell[{{p{\linewidth}}}]{
                \texttt{\tiny{[GM$|$GM]}}
                \texttt{حرّك: C, (4,7)} \\
            }
        }
    }
    & & \\ \\

    \theutterance \stepcounter{utterance}  
    & & & \multicolumn{2}{p{0.3\linewidth}}{
        \cellcolor[rgb]{0.9,0.9,0.9}{
            \makecell[{{p{\linewidth}}}]{
                \texttt{\tiny{[GM$|$GM]}}
                \texttt{C moved to (4, 7)} \\
            }
        }
    }
    & & \\ \\

    \theutterance \stepcounter{utterance}  
    & & \multicolumn{4}{p{0.6\linewidth}}{
        \cellcolor[rgb]{0.9,0.9,0.9}{
            \makecell[{{p{\linewidth}}}]{
                \texttt{\tiny{[P1$\langle$GM]}}
                \texttt{تم تحريك 'C' إلى (7,5) بنجاح. تبدو شبكتك المحدّثة هكذا:} \\
\texttt{```} \\
\texttt{ 1234567} \\
\texttt{╔═╤══╤══╗} \\
\texttt{║◌│◌◌│◌◌║ 1} \\
\texttt{╟─┤◌◌│◌◌║ 2} \\
\texttt{║◌│◌◌├──╢ 3} \\
\texttt{╟─┘◌◌│◌◌║ 4} \\
\texttt{║◌◌◌◌│◌C║ 5} \\
\texttt{║◌L┌─┼──╢ 6} \\
\texttt{║◌◌│◌│P◌║ 7} \\
\texttt{╚══╧═╧══╝} \\
\texttt{```} \\
\texttt{أنت تلعب حالياً الجولة 4 من أصل 12 كحد أقصى.} \\
\texttt{لقد جمعتما معاً 5 من أصل 8 عقوبة.} \\
\texttt{حرّك اللاعب الآخر كائناً على شبكته.} \\
\\ 
\\ 
\texttt{ما هو أمرك التالي؟} \\
            }
        }
    }
    & \\ \\

    \theutterance \stepcounter{utterance}  
    & \multicolumn{4}{p{0.6\linewidth}}{
        \cellcolor[rgb]{0.9,0.9,0.9}{
            \makecell[{{p{\linewidth}}}]{
                \texttt{\tiny{[P1$\rangle$GM]}}
                \texttt{حرّك: P, (7,7)} \\
            }
        }
    }
    & & \\ \\

    \theutterance \stepcounter{utterance}  
    & & & \multicolumn{2}{p{0.3\linewidth}}{
        \cellcolor[rgb]{0.9,0.9,0.9}{
            \makecell[{{p{\linewidth}}}]{
                \texttt{\tiny{[GM$|$GM]}}
                \texttt{حرّك: P, (7,7)} \\
            }
        }
    }
    & & \\ \\

    \theutterance \stepcounter{utterance}  
    & & & \multicolumn{2}{p{0.3\linewidth}}{
        \cellcolor[rgb]{0.9,0.9,0.9}{
            \makecell[{{p{\linewidth}}}]{
                \texttt{\tiny{[GM$|$GM]}}
                \texttt{P moved to (7, 7)} \\
            }
        }
    }
    & & \\ \\

    \theutterance \stepcounter{utterance}  
    & & \multicolumn{4}{p{0.6\linewidth}}{
        \cellcolor[rgb]{0.9,0.9,0.9}{
            \makecell[{{p{\linewidth}}}]{
                \texttt{\tiny{[GM$\rangle$P2]}}
                \texttt{تم تحريك 'C' إلى (4,7) بنجاح. تبدو شبكتك المحدّثة هكذا:} \\
\texttt{```} \\
\texttt{ 1234567} \\
\texttt{╔═╤══╤══╗} \\
\texttt{║◌│◌◌│◌◌║ 1} \\
\texttt{╟─┤◌◌│◌◌║ 2} \\
\texttt{║◌│◌◌├──╢ 3} \\
\texttt{╟─┘◌◌│◌◌║ 4} \\
\texttt{║P◌◌◌│◌◌║ 5} \\
\texttt{║◌◌┌─┼──╢ 6} \\
\texttt{║◌◌│C│◌L║ 7} \\
\texttt{╚══╧═╧══╝} \\
\texttt{```} \\
\texttt{أنت تلعب حالياً الجولة 5 من أصل 12 كحد أقصى.} \\
\texttt{لقد جمعتما معاً 5 من أصل 8 عقوبة.} \\
\texttt{حرّك اللاعب الآخر كائناً على شبكته.} \\
\\ 
\\ 
\texttt{ما هو أمرك التالي؟} \\
            }
        }
    }
    & \\ \\

    \theutterance \stepcounter{utterance}  
    & & & \multicolumn{4}{p{0.6\linewidth}}{
        \cellcolor[rgb]{0.9,0.9,0.9}{
            \makecell[{{p{\linewidth}}}]{
                \texttt{\tiny{[GM$\langle$P2]}}
                \texttt{حرّك: P, (7,7)} \\
            }
        }
    }
     \\ \\

    \theutterance \stepcounter{utterance}  
    & & & \multicolumn{2}{p{0.3\linewidth}}{
        \cellcolor[rgb]{0.9,0.9,0.9}{
            \makecell[{{p{\linewidth}}}]{
                \texttt{\tiny{[GM$|$GM]}}
                \texttt{حرّك: P, (7,7)} \\
            }
        }
    }
    & & \\ \\

    \theutterance \stepcounter{utterance}  
    & & & \multicolumn{2}{p{0.3\linewidth}}{
        \cellcolor[rgb]{0.9,0.9,0.9}{
            \makecell[{{p{\linewidth}}}]{
                \texttt{\tiny{[GM$|$GM]}}
                \texttt{عقوبة: (7,7) ليست فارغة، بل تحتوي على 'L'.} \\
\\ 
\texttt{لقد جمعتما معاً 6 من أصل 8 عقوبة.} \\
\texttt{يرجى المحاولة مرة أخرى!} \\
            }
        }
    }
    & & \\ \\

    \theutterance \stepcounter{utterance}  
    & & \multicolumn{4}{p{0.6\linewidth}}{
        \cellcolor[rgb]{0.9,0.9,0.9}{
            \makecell[{{p{\linewidth}}}]{
                \texttt{\tiny{[GM$\rangle$P2]}}
                \texttt{عقوبة: (7,7) ليست فارغة، بل تحتوي على 'L'.} \\
\\ 
\texttt{لقد جمعتما معاً 6 من أصل 8 عقوبة.} \\
\texttt{يرجى المحاولة مرة أخرى!} \\
            }
        }
    }
    & \\ \\

    \theutterance \stepcounter{utterance}  
    & & & \multicolumn{4}{p{0.6\linewidth}}{
        \cellcolor[rgb]{0.9,0.9,0.9}{
            \makecell[{{p{\linewidth}}}]{
                \texttt{\tiny{[GM$\langle$P2]}}
                \texttt{حرّك: P, (7,1)} \\
            }
        }
    }
     \\ \\

    \theutterance \stepcounter{utterance}  
    & & & \multicolumn{2}{p{0.3\linewidth}}{
        \cellcolor[rgb]{0.9,0.9,0.9}{
            \makecell[{{p{\linewidth}}}]{
                \texttt{\tiny{[GM$|$GM]}}
                \texttt{حرّك: P, (7,1)} \\
            }
        }
    }
    & & \\ \\

    \theutterance \stepcounter{utterance}  
    & & & \multicolumn{2}{p{0.3\linewidth}}{
        \cellcolor[rgb]{0.9,0.9,0.9}{
            \makecell[{{p{\linewidth}}}]{
                \texttt{\tiny{[GM$|$GM]}}
                \texttt{P moved to (7, 1)} \\
            }
        }
    }
    & & \\ \\

    \theutterance \stepcounter{utterance}  
    & & \multicolumn{4}{p{0.6\linewidth}}{
        \cellcolor[rgb]{0.9,0.9,0.9}{
            \makecell[{{p{\linewidth}}}]{
                \texttt{\tiny{[P1$\langle$GM]}}
                \texttt{تم تحريك 'P' إلى (7,7) بنجاح. تبدو شبكتك المحدّثة هكذا:} \\
\texttt{```} \\
\texttt{ 1234567} \\
\texttt{╔═╤══╤══╗} \\
\texttt{║◌│◌◌│◌◌║ 1} \\
\texttt{╟─┤◌◌│◌◌║ 2} \\
\texttt{║◌│◌◌├──╢ 3} \\
\texttt{╟─┘◌◌│◌◌║ 4} \\
\texttt{║◌◌◌◌│◌C║ 5} \\
\texttt{║◌L┌─┼──╢ 6} \\
\texttt{║◌◌│◌│◌P║ 7} \\
\texttt{╚══╧═╧══╝} \\
\texttt{```} \\
\texttt{أنت تلعب حالياً الجولة 5 من أصل 12 كحد أقصى.} \\
\texttt{لقد جمعتما معاً 6 من أصل 8 عقوبة.} \\
\texttt{حرّك اللاعب الآخر كائناً على شبكته.} \\
\\ 
\\ 
\texttt{ما هو أمرك التالي؟} \\
            }
        }
    }
    & \\ \\

    \theutterance \stepcounter{utterance}  
    & \multicolumn{4}{p{0.6\linewidth}}{
        \cellcolor[rgb]{0.9,0.9,0.9}{
            \makecell[{{p{\linewidth}}}]{
                \texttt{\tiny{[P1$\rangle$GM]}}
                \texttt{حرّك: C, (7,6)} \\
            }
        }
    }
    & & \\ \\

    \theutterance \stepcounter{utterance}  
    & & & \multicolumn{2}{p{0.3\linewidth}}{
        \cellcolor[rgb]{0.9,0.9,0.9}{
            \makecell[{{p{\linewidth}}}]{
                \texttt{\tiny{[GM$|$GM]}}
                \texttt{حرّك: C, (7,6)} \\
            }
        }
    }
    & & \\ \\

    \theutterance \stepcounter{utterance}  
    & & & \multicolumn{2}{p{0.3\linewidth}}{
        \cellcolor[rgb]{0.9,0.9,0.9}{
            \makecell[{{p{\linewidth}}}]{
                \texttt{\tiny{[GM$|$GM]}}
                \texttt{عقوبة: (7,6) ليست فارغة، بل تحتوي على '─'.} \\
\\ 
\texttt{لقد جمعتما معاً 7 من أصل 8 عقوبة.} \\
\texttt{يرجى المحاولة مرة أخرى!} \\
            }
        }
    }
    & & \\ \\

    \theutterance \stepcounter{utterance}  
    & & \multicolumn{4}{p{0.6\linewidth}}{
        \cellcolor[rgb]{0.9,0.9,0.9}{
            \makecell[{{p{\linewidth}}}]{
                \texttt{\tiny{[P1$\langle$GM]}}
                \texttt{عقوبة: (7,6) ليست فارغة، بل تحتوي على '─'.} \\
\\ 
\texttt{لقد جمعتما معاً 7 من أصل 8 عقوبة.} \\
\texttt{يرجى المحاولة مرة أخرى!} \\
            }
        }
    }
    & \\ \\

    \theutterance \stepcounter{utterance}  
    & \multicolumn{4}{p{0.6\linewidth}}{
        \cellcolor[rgb]{0.9,0.9,0.9}{
            \makecell[{{p{\linewidth}}}]{
                \texttt{\tiny{[P1$\rangle$GM]}}
                \texttt{حرّك: C, (7,4)} \\
            }
        }
    }
    & & \\ \\

    \theutterance \stepcounter{utterance}  
    & & & \multicolumn{2}{p{0.3\linewidth}}{
        \cellcolor[rgb]{0.9,0.9,0.9}{
            \makecell[{{p{\linewidth}}}]{
                \texttt{\tiny{[GM$|$GM]}}
                \texttt{حرّك: C, (7,4)} \\
            }
        }
    }
    & & \\ \\

    \theutterance \stepcounter{utterance}  
    & & & \multicolumn{2}{p{0.3\linewidth}}{
        \cellcolor[rgb]{0.9,0.9,0.9}{
            \makecell[{{p{\linewidth}}}]{
                \texttt{\tiny{[GM$|$GM]}}
                \texttt{C moved to (7, 4)} \\
            }
        }
    }
    & & \\ \\

    \theutterance \stepcounter{utterance}  
    & & \multicolumn{4}{p{0.6\linewidth}}{
        \cellcolor[rgb]{0.9,0.9,0.9}{
            \makecell[{{p{\linewidth}}}]{
                \texttt{\tiny{[GM$\rangle$P2]}}
                \texttt{تم تحريك 'P' إلى (7,1) بنجاح. تبدو شبكتك المحدّثة هكذا:} \\
\texttt{```} \\
\texttt{ 1234567} \\
\texttt{╔═╤══╤══╗} \\
\texttt{║◌│◌◌│◌P║ 1} \\
\texttt{╟─┤◌◌│◌◌║ 2} \\
\texttt{║◌│◌◌├──╢ 3} \\
\texttt{╟─┘◌◌│◌◌║ 4} \\
\texttt{║◌◌◌◌│◌◌║ 5} \\
\texttt{║◌◌┌─┼──╢ 6} \\
\texttt{║◌◌│C│◌L║ 7} \\
\texttt{╚══╧═╧══╝} \\
\texttt{```} \\
\texttt{أنت تلعب حالياً الجولة 6 من أصل 12 كحد أقصى.} \\
\texttt{لقد جمعتما معاً 7 من أصل 8 عقوبة.} \\
\texttt{حرّك اللاعب الآخر كائناً على شبكته.} \\
\\ 
\\ 
\texttt{ما هو أمرك التالي؟} \\
            }
        }
    }
    & \\ \\

    \theutterance \stepcounter{utterance}  
    & & & \multicolumn{4}{p{0.6\linewidth}}{
        \cellcolor[rgb]{0.9,0.9,0.9}{
            \makecell[{{p{\linewidth}}}]{
                \texttt{\tiny{[GM$\langle$P2]}}
                \texttt{قل: انتهيت؟} \\
            }
        }
    }
     \\ \\

    \theutterance \stepcounter{utterance}  
    & & & \multicolumn{2}{p{0.3\linewidth}}{
        \cellcolor[rgb]{0.9,0.9,0.9}{
            \makecell[{{p{\linewidth}}}]{
                \texttt{\tiny{[GM$|$GM]}}
                \texttt{قل: انتهيت؟} \\
            }
        }
    }
    & & \\ \\

    \theutterance \stepcounter{utterance}  
    & & \multicolumn{4}{p{0.6\linewidth}}{
        \cellcolor[rgb]{0.9,0.9,0.9}{
            \makecell[{{p{\linewidth}}}]{
                \texttt{\tiny{[P1$\langle$GM]}}
                \texttt{تم تحريك 'C' إلى (7,4) بنجاح. تبدو شبكتك المحدّثة هكذا:} \\
\texttt{```} \\
\texttt{ 1234567} \\
\texttt{╔═╤══╤══╗} \\
\texttt{║◌│◌◌│◌◌║ 1} \\
\texttt{╟─┤◌◌│◌◌║ 2} \\
\texttt{║◌│◌◌├──╢ 3} \\
\texttt{╟─┘◌◌│◌C║ 4} \\
\texttt{║◌◌◌◌│◌◌║ 5} \\
\texttt{║◌L┌─┼──╢ 6} \\
\texttt{║◌◌│◌│◌P║ 7} \\
\texttt{╚══╧═╧══╝} \\
\texttt{```} \\
\texttt{أنت تلعب حالياً الجولة 6 من أصل 12 كحد أقصى.} \\
\texttt{لقد جمعتما معاً 7 من أصل 8 عقوبة.} \\
\texttt{أرسل اللاعب الآخر هذه الرسالة:} \\
\\ 
\\ 
\texttt{"انتهيت؟"} \\
\\ 
\\ 
\texttt{ما هو أمرك التالي؟} \\
            }
        }
    }
    & \\ \\

    \theutterance \stepcounter{utterance}  
    & \multicolumn{4}{p{0.6\linewidth}}{
        \cellcolor[rgb]{0.9,0.9,0.9}{
            \makecell[{{p{\linewidth}}}]{
                \texttt{\tiny{[P1$\rangle$GM]}}
                \texttt{قل: انتهيت!} \\
            }
        }
    }
    & & \\ \\

    \theutterance \stepcounter{utterance}  
    & & & \multicolumn{2}{p{0.3\linewidth}}{
        \cellcolor[rgb]{0.9,0.9,0.9}{
            \makecell[{{p{\linewidth}}}]{
                \texttt{\tiny{[GM$|$GM]}}
                \texttt{قل: انتهيت!} \\
            }
        }
    }
    & & \\ \\

    \theutterance \stepcounter{utterance}  
    & & & \multicolumn{2}{p{0.3\linewidth}}{
        \cellcolor[rgb]{0.9,0.9,0.9}{
            \makecell[{{p{\linewidth}}}]{
                \texttt{\tiny{[GM$|$GM]}}
                \texttt{true} \\
            }
        }
    }
    & & \\ \\

    \theutterance \stepcounter{utterance}  
    & & & \multicolumn{2}{p{0.3\linewidth}}{
        \cellcolor[rgb]{0.9,0.9,0.9}{
            \makecell[{{p{\linewidth}}}]{
                \texttt{\tiny{[GM$|$GM]}}
                \texttt{Initial states:} \\
\texttt{```} \\
\texttt{Player 1:} \\
\texttt{ 1234567} \\
\texttt{╔═╤══╤══╗} \\
\texttt{║◌│◌◌│◌◌║ 1} \\
\texttt{╟─┤◌◌│◌◌║ 2} \\
\texttt{║◌│◌◌├──╢ 3} \\
\texttt{╟─┘C◌│◌◌║ 4} \\
\texttt{║◌◌◌◌│◌◌║ 5} \\
\texttt{║◌L┌─┼──╢ 6} \\
\texttt{║◌◌│◌│P◌║ 7} \\
\texttt{╚══╧═╧══╝} \\
\\ 
\\ 
\texttt{Player 2:} \\
\texttt{ 1234567} \\
\texttt{╔═╤══╤══╗} \\
\texttt{║◌│◌◌│◌◌║ 1} \\
\texttt{╟─┤◌◌│◌◌║ 2} \\
\texttt{║◌│L◌├──╢ 3} \\
\texttt{╟─┘◌◌│◌C║ 4} \\
\texttt{║P◌◌◌│◌◌║ 5} \\
\texttt{║◌◌┌─┼──╢ 6} \\
\texttt{║◌◌│◌│◌◌║ 7} \\
\texttt{╚══╧═╧══╝} \\
\\ 
\texttt{```} \\
            }
        }
    }
    & & \\ \\

    \theutterance \stepcounter{utterance}  
    & & & \multicolumn{2}{p{0.3\linewidth}}{
        \cellcolor[rgb]{0.9,0.9,0.9}{
            \makecell[{{p{\linewidth}}}]{
                \texttt{\tiny{[GM$|$GM]}}
                \texttt{End states:} \\
\texttt{```} \\
\texttt{Player 1:} \\
\texttt{ 1234567} \\
\texttt{╔═╤══╤══╗} \\
\texttt{║◌│◌◌│◌◌║ 1} \\
\texttt{╟─┤◌◌│◌◌║ 2} \\
\texttt{║◌│◌◌├──╢ 3} \\
\texttt{╟─┘◌◌│◌C║ 4} \\
\texttt{║◌◌◌◌│◌◌║ 5} \\
\texttt{║◌L┌─┼──╢ 6} \\
\texttt{║◌◌│◌│◌P║ 7} \\
\texttt{╚══╧═╧══╝} \\
\\ 
\\ 
\texttt{Player 2:} \\
\texttt{ 1234567} \\
\texttt{╔═╤══╤══╗} \\
\texttt{║◌│◌◌│◌P║ 1} \\
\texttt{╟─┤◌◌│◌◌║ 2} \\
\texttt{║◌│◌◌├──╢ 3} \\
\texttt{╟─┘◌◌│◌◌║ 4} \\
\texttt{║◌◌◌◌│◌◌║ 5} \\
\texttt{║◌◌┌─┼──╢ 6} \\
\texttt{║◌◌│C│◌L║ 7} \\
\texttt{╚══╧═╧══╝} \\
\\ 
\texttt{```} \\
            }
        }
    }
    & & \\ \\

    \theutterance \stepcounter{utterance}  
    & & & \multicolumn{2}{p{0.3\linewidth}}{
        \cellcolor[rgb]{0.9,0.9,0.9}{
            \makecell[{{p{\linewidth}}}]{
                \texttt{\tiny{[GM$|$GM]}}
                \texttt{* success: True} \\
\texttt{* lose: True} \\
\texttt{* aborted: False} \\
\texttt{{-}{-}{-}{-}{-}{-}{-}} \\
\texttt{* Shifts: 4.00} \\
\texttt{* Max Shifts: 4.00} \\
\texttt{* Min Shifts: 2.00} \\
\texttt{* End Distance Sum: 15.34} \\
\texttt{* Init Distance Sum: 12.55} \\
\texttt{* Expected Distance Sum: 12.57} \\
\texttt{* Penalties: 7.00} \\
\texttt{* Max Penalties: 8.00} \\
\texttt{* Rounds: 6.00} \\
\texttt{* Max Rounds: 12.00} \\
\texttt{* Object Count: 3.00} \\
            }
        }
    }
    & & \\ \\

\end{supertabular}
}

\end{document}
